\documentclass[english,screen,final]{noahassignment}

\title{Assignment 2: Mechanical Vibrations}
\author{Example Student}
\studentid{202600001}

\documentsetup{
  course={Mechanical Vibrations},
  course-code={MECH-402},
  course-short={Vibrations},
  running-title={Assignment 2},
  instructor={Grace Hopper},
  semester={Autumn 2026},
  submission-date={15 October 2026},
  submission-date-iso={2026-10-15},
  logo={../../figures/AU.pdf},
  problem-numbering=automatic,
  subproblem-numbering=letters,
  toc-subproblems=true
}

\begin{document}

\frontmatter
\maketitle
\tableofcontents

\mainmatter

\lecture[Vibrations]{1}{8 October 2026}{Single-degree-of-freedom vibration}

\begin{problem}\label{prob:oscillator}
A mass \(m\) is connected to a linear spring with stiffness \(k\).
Neglect damping.

\begin{subproblems}
  \subproblem Derive the governing equation.\label{prob:oscillator:equation}
  \subproblem Determine the natural frequency.\label{prob:oscillator:frequency}
\end{subproblems}
\end{problem}

\begin{assumptions}
  \item The spring is linear.
  \item Motion is one-dimensional.
\end{assumptions}

\begin{givens}
  m & \qty{2}{\kilogram} \\
  k & \qty{18}{\newton\per\meter}
\end{givens}

\begin{derivation}
  \derivationpart[prob:oscillator:equation]
  Newton's second law gives
  \[
    m\ddot{x}+kx=0.
  \]
  \begin{subresult}
    \[
      m\ddot{x}+kx=0
    \]
  \end{subresult}

  \derivationpart[prob:oscillator:frequency]
  Comparing with the standard oscillator equation gives
  \[
    \omega_n^2=\frac{k}{m}.
  \]
  \begin{subresult}
    \[
      \omega_n=\sqrt{\frac{k}{m}}=\qty{3}{\radian\per\second}
    \]
  \end{subresult}
\end{derivation}

\end{document}
